\subsection{Nonuniqueness in the minimal envelope class}

In this subsection the manifold and the background are no longer fixed. The examples of \cref{uni-thm:analytic-example} are smooth solutions on projective manifolds with nonclosed Hermitian backgrounds. Comparison with a K\"ahler form shows that these backgrounds have both positive lower mass and bounded mass, and smooth potentials belong to the minimal envelope class. Thus even both mass bounds together allow several normalized solutions with the same scalar.

\begin{corollary}\label{cor:analytic-nonuniqueness-envelope-class}
In every complex dimension $n\geq3$, there are a projective manifold $X$ of dimension $n$, a smooth nonclosed Hermitian form $\omega$ on $X$, and a nonzero finite non-pluripolar measure $\nu$ on $X$ such that
\[
 0<v_-(\omega)\leq v_+(\omega)<+\infty,
\]
but the equation
\[
 \mathrm{MA}^{\mathrm{ext}}_\omega(u)=c\nu,
 \qquad u\in\mathcal E_{\mathrm{env}}(X,\omega,0),
 \qquad \sup_Xu=0,
\]
has at least two smooth solutions $u$ for $c=1$, and no solution for any other positive constant $c$. Relative to a K\"ahler volume form on $X$, the measure $\nu$ has a real analytic density, positive off two connected smooth real hypersurfaces and vanishing quadratically along them.
\end{corollary}
\begin{proof}
Use the form and the two potentials of \cref{uni-thm:analytic-example}. Since the manifold is projective, \cref{lem:kahler-modification-bounded-mass} gives bounded mass. Fix a K\"ahler form $\Omega$ on $X$ and choose $\varepsilon>0$ with $\varepsilon\Omega\leq\omega$. Lower-volume monotonicity, \cref{thm:volume-monotonicity}, applied to the zero potential in both backgrounds, gives
\[
 v_-(\omega)\geq v_-(\varepsilon\Omega)
 =\varepsilon^n\int_X\Omega^n>0.
\]
Here the potentials of the same ordinary type as zero are exactly the bounded ones, so these lower volumes are the infima defining $v_-$.

The two smooth potentials belong to $\mathcal E_{\mathrm{env}}(X,\omega,0)$ by \cref{prop:relative-elementary}, and their extended measures are the ordinary top powers. Their common measure $\nu$ is finite and does not charge pluripolar sets. The construction gives $c=1$, and \cref{cor:relative-extended-exponential-order} shows that no other positive scalar is possible. The density and its zero set were computed in \cref{uni-thm:analytic-example}.
\end{proof}
